% Studies-in-Algebra-and-Group-Theory - Lecture 6: Kernals, Images, and more
% Author: S. D. V. Stephens
% Date: 2025-12-21
% Compile with: pdflatex -shell-escape

\documentclass{report}
\input{~/university/preamble.tex}
% preamble-darkmode.tex
% Dark mode extension for lecture notes
% Usage: % preamble-darkmode.tex
% Dark mode extension for lecture notes
% Usage: % preamble-darkmode.tex
% Dark mode extension for lecture notes
% Usage: \input{~/university/preamble-darkmode.tex} AFTER preamble.tex
%
% This provides:
% - Dark background with light text
% - Material Design color palette
% - Ocean blue gradient colors (p1-p9)
% - Enhanced TikZ/PGFPlots for dark backgrounds
% - qq tcolorbox environment
% - tkz-euclide for geometric constructions

% ===========================================
% DARK MODE PACKAGE
% ===========================================
\usepackage{darkmode}
\enabledarkmode

% ===========================================
% ADDITIONAL PACKAGES
% ===========================================
\usepackage{changepage}
\usepackage{ulem}
\usepackage{colortbl}
\usepackage{tkz-euclide}
\usepackage{capt-of}

% ===========================================
% EXTENDED TIKZ LIBRARIES
% ===========================================
\usetikzlibrary{patterns,3d,backgrounds}
\usepgfplotslibrary{groupplots}

% Upgrade pgfplots compatibility
\pgfplotsset{compat=1.18}

% ===========================================
% OCEAN BLUE GRADIENT (p1-p9)
% ===========================================
\definecolor{p1}{HTML}{caf0f8}
\definecolor{p2}{HTML}{ade8f4}
\definecolor{p3}{HTML}{90e0ef}
\definecolor{p4}{HTML}{48cae4}
\definecolor{p5}{HTML}{00b4d8}
\definecolor{p6}{HTML}{0096c7}
\definecolor{p7}{HTML}{0077b6}
\definecolor{p8}{HTML}{023e8a}
\definecolor{p9}{HTML}{03045e}

% Dark background color
\definecolor{pag}{HTML}{293133}

% ===========================================
% MATERIAL DESIGN COLORS
% ===========================================
% Note: myred, myyellow already defined in main preamble
% These are additional/alternate versions
\definecolor{matred}{HTML}{f44336}
\definecolor{matpink}{HTML}{e81e63}
\definecolor{matpurple}{HTML}{9c27b0}
\definecolor{matdeeppurple}{HTML}{673ab7}
\definecolor{matindigo}{HTML}{3f51b5}
\definecolor{matblue}{HTML}{2196f3}
\definecolor{matlightblue}{HTML}{03a9f4}
\definecolor{matcyan}{HTML}{00bcd4}
\definecolor{matteal}{HTML}{009688}
\definecolor{matgreen}{HTML}{4caf50}
\definecolor{matlightgreen}{HTML}{8bc34a}
\definecolor{matlime}{HTML}{cddc39}
\definecolor{matyellow}{HTML}{ffeb3b}
\definecolor{matamber}{HTML}{ffc107}
\definecolor{matorange}{HTML}{ff9800}
\definecolor{matdeeporange}{HTML}{ff5722}
\definecolor{matbrown}{HTML}{795548}
\definecolor{matgray}{HTML}{9e9e9e}
\definecolor{matbluegray}{HTML}{607d8b}

% ===========================================
% COLORLET FOR TIKZ
% ===========================================
\colorlet{xcol}{blue!60!black}

% ===========================================
% DARK MODE TCOLORBOX
% ===========================================
\newtcolorbox{qq}[2][]{%
    boxrule=0.75pt,
    sharp corners,
    colframe=white,
    colback=pag,
    coltext=white,
    #1,
}

% Dark definition box
\newtcolorbox{darkdfn}[2][]{%
    enhanced,
    boxrule=1pt,
    sharp corners,
    colframe=p5,
    colback=pag,
    coltext=white,
    coltitle=white,
    fonttitle=\bfseries,
    title=#2,
    #1,
}

% Dark theorem box
\newtcolorbox{darkthm}[2][]{%
    enhanced,
    boxrule=1pt,
    sharp corners,
    colframe=matpurple,
    colback=pag,
    coltext=white,
    coltitle=white,
    fonttitle=\bfseries,
    title=#2,
    #1,
}

% ===========================================
% TIKZ 3D FIX
% ===========================================
% Small fix for canvas is xy plane at z
% https://tex.stackexchange.com/a/48776/121799
\makeatletter
\tikzoption{canvas is xy plane at z}[]{%
    \def\tikz@plane@origin{\pgfpointxyz{0}{0}{#1}}%
    \def\tikz@plane@x{\pgfpointxyz{1}{0}{#1}}%
    \def\tikz@plane@y{\pgfpointxyz{0}{1}{#1}}%
    \tikz@canvas@is@plane}
\makeatother

% ===========================================
% CONDITIONAL LINE TIKZSET
% ===========================================
\tikzset{
    conditional line/.style 2 args={
        /utils/exec={\pgfmathsetmacro{\xcoordA}{#1}},
        /utils/exec={\pgfmathsetmacro{\xcoordB}{#2}},
        insert path={
            \ifdim\xcoordA pt>\xcoordB pt
            (\xcoordA,0) -- (\xcoordB,0)
            \fi
        },
    },
}

% ===========================================
% PGFPLOTS DARK MODE DEFAULTS
% ===========================================
\pgfplotsset{
    dark axis/.style={
        axis line style={white},
        tick style={white},
        label style={white},
        grid style={pag!70},
        every axis label/.append style={white},
        every tick label/.append style={white},
    },
}

% ===========================================
% TIKZ EXTERNALIZATION (for fast compilation)
% ===========================================
% This creates .md5 cache files for figures
% Requires: pdflatex -shell-escape
%
% To enable, add AFTER loading this file:
%   \enabletikzexternalize
%
\newcommand{\enabletikzexternalize}{%
  \usetikzlibrary{external}%
  \tikzexternalize[prefix=figures/]%
}

% ===========================================
% CONVENIENT SHORTCUTS FOR DARK PLOTS
% ===========================================
\newcommand{\darkplot}[1]{%
    \begin{tikzpicture}
    \begin{axis}[
        dark axis,
        grid=both,
        major grid style={line width=.2pt,draw=pag!70},
        #1
    ]
}


% ===========================================
% QUICK GEOMETRY MACROS (tkz-euclide)
% ===========================================
% Draw a point: \point{name}{x}{y}
\newcommand{\gpoint}[3]{\tkzDefPoint(#2,#3){#1}}

% Draw line through points: \gline{A}{B}
\newcommand{\gline}[2]{\tkzDrawLine[color=p4](#1,#2)}

% Draw circle: \gcircle{center}{radius}
\newcommand{\gcircle}[2]{\tkzDrawCircle[color=p5](#1,#2)}

% Mark angle: \gangle{A}{B}{C}
\newcommand{\gangle}[3]{\tkzMarkAngle[color=p3,size=0.5](#1,#2,#3)}

% ===========================================
% USAGE EXAMPLE
% ===========================================
% In your document:
%
% \begin{tikzpicture}
%   \begin{axis}[dark axis, ...]
%     \addplot[p4, thick] {sin(deg(x))};
%   \end{axis}
% \end{tikzpicture}
%
% Or with qq box:
% \begin{qq}{My Title}
%   Content here in dark box
% \end{qq}
 AFTER preamble.tex
%
% This provides:
% - Dark background with light text
% - Material Design color palette
% - Ocean blue gradient colors (p1-p9)
% - Enhanced TikZ/PGFPlots for dark backgrounds
% - qq tcolorbox environment
% - tkz-euclide for geometric constructions

% ===========================================
% DARK MODE PACKAGE
% ===========================================
\usepackage{darkmode}
\enabledarkmode

% ===========================================
% ADDITIONAL PACKAGES
% ===========================================
\usepackage{changepage}
\usepackage{ulem}
\usepackage{colortbl}
\usepackage{tkz-euclide}
\usepackage{capt-of}

% ===========================================
% EXTENDED TIKZ LIBRARIES
% ===========================================
\usetikzlibrary{patterns,3d,backgrounds}
\usepgfplotslibrary{groupplots}

% Upgrade pgfplots compatibility
\pgfplotsset{compat=1.18}

% ===========================================
% OCEAN BLUE GRADIENT (p1-p9)
% ===========================================
\definecolor{p1}{HTML}{caf0f8}
\definecolor{p2}{HTML}{ade8f4}
\definecolor{p3}{HTML}{90e0ef}
\definecolor{p4}{HTML}{48cae4}
\definecolor{p5}{HTML}{00b4d8}
\definecolor{p6}{HTML}{0096c7}
\definecolor{p7}{HTML}{0077b6}
\definecolor{p8}{HTML}{023e8a}
\definecolor{p9}{HTML}{03045e}

% Dark background color
\definecolor{pag}{HTML}{293133}

% ===========================================
% MATERIAL DESIGN COLORS
% ===========================================
% Note: myred, myyellow already defined in main preamble
% These are additional/alternate versions
\definecolor{matred}{HTML}{f44336}
\definecolor{matpink}{HTML}{e81e63}
\definecolor{matpurple}{HTML}{9c27b0}
\definecolor{matdeeppurple}{HTML}{673ab7}
\definecolor{matindigo}{HTML}{3f51b5}
\definecolor{matblue}{HTML}{2196f3}
\definecolor{matlightblue}{HTML}{03a9f4}
\definecolor{matcyan}{HTML}{00bcd4}
\definecolor{matteal}{HTML}{009688}
\definecolor{matgreen}{HTML}{4caf50}
\definecolor{matlightgreen}{HTML}{8bc34a}
\definecolor{matlime}{HTML}{cddc39}
\definecolor{matyellow}{HTML}{ffeb3b}
\definecolor{matamber}{HTML}{ffc107}
\definecolor{matorange}{HTML}{ff9800}
\definecolor{matdeeporange}{HTML}{ff5722}
\definecolor{matbrown}{HTML}{795548}
\definecolor{matgray}{HTML}{9e9e9e}
\definecolor{matbluegray}{HTML}{607d8b}

% ===========================================
% COLORLET FOR TIKZ
% ===========================================
\colorlet{xcol}{blue!60!black}

% ===========================================
% DARK MODE TCOLORBOX
% ===========================================
\newtcolorbox{qq}[2][]{%
    boxrule=0.75pt,
    sharp corners,
    colframe=white,
    colback=pag,
    coltext=white,
    #1,
}

% Dark definition box
\newtcolorbox{darkdfn}[2][]{%
    enhanced,
    boxrule=1pt,
    sharp corners,
    colframe=p5,
    colback=pag,
    coltext=white,
    coltitle=white,
    fonttitle=\bfseries,
    title=#2,
    #1,
}

% Dark theorem box
\newtcolorbox{darkthm}[2][]{%
    enhanced,
    boxrule=1pt,
    sharp corners,
    colframe=matpurple,
    colback=pag,
    coltext=white,
    coltitle=white,
    fonttitle=\bfseries,
    title=#2,
    #1,
}

% ===========================================
% TIKZ 3D FIX
% ===========================================
% Small fix for canvas is xy plane at z
% https://tex.stackexchange.com/a/48776/121799
\makeatletter
\tikzoption{canvas is xy plane at z}[]{%
    \def\tikz@plane@origin{\pgfpointxyz{0}{0}{#1}}%
    \def\tikz@plane@x{\pgfpointxyz{1}{0}{#1}}%
    \def\tikz@plane@y{\pgfpointxyz{0}{1}{#1}}%
    \tikz@canvas@is@plane}
\makeatother

% ===========================================
% CONDITIONAL LINE TIKZSET
% ===========================================
\tikzset{
    conditional line/.style 2 args={
        /utils/exec={\pgfmathsetmacro{\xcoordA}{#1}},
        /utils/exec={\pgfmathsetmacro{\xcoordB}{#2}},
        insert path={
            \ifdim\xcoordA pt>\xcoordB pt
            (\xcoordA,0) -- (\xcoordB,0)
            \fi
        },
    },
}

% ===========================================
% PGFPLOTS DARK MODE DEFAULTS
% ===========================================
\pgfplotsset{
    dark axis/.style={
        axis line style={white},
        tick style={white},
        label style={white},
        grid style={pag!70},
        every axis label/.append style={white},
        every tick label/.append style={white},
    },
}

% ===========================================
% TIKZ EXTERNALIZATION (for fast compilation)
% ===========================================
% This creates .md5 cache files for figures
% Requires: pdflatex -shell-escape
%
% To enable, add AFTER loading this file:
%   \enabletikzexternalize
%
\newcommand{\enabletikzexternalize}{%
  \usetikzlibrary{external}%
  \tikzexternalize[prefix=figures/]%
}

% ===========================================
% CONVENIENT SHORTCUTS FOR DARK PLOTS
% ===========================================
\newcommand{\darkplot}[1]{%
    \begin{tikzpicture}
    \begin{axis}[
        dark axis,
        grid=both,
        major grid style={line width=.2pt,draw=pag!70},
        #1
    ]
}


% ===========================================
% QUICK GEOMETRY MACROS (tkz-euclide)
% ===========================================
% Draw a point: \point{name}{x}{y}
\newcommand{\gpoint}[3]{\tkzDefPoint(#2,#3){#1}}

% Draw line through points: \gline{A}{B}
\newcommand{\gline}[2]{\tkzDrawLine[color=p4](#1,#2)}

% Draw circle: \gcircle{center}{radius}
\newcommand{\gcircle}[2]{\tkzDrawCircle[color=p5](#1,#2)}

% Mark angle: \gangle{A}{B}{C}
\newcommand{\gangle}[3]{\tkzMarkAngle[color=p3,size=0.5](#1,#2,#3)}

% ===========================================
% USAGE EXAMPLE
% ===========================================
% In your document:
%
% \begin{tikzpicture}
%   \begin{axis}[dark axis, ...]
%     \addplot[p4, thick] {sin(deg(x))};
%   \end{axis}
% \end{tikzpicture}
%
% Or with qq box:
% \begin{qq}{My Title}
%   Content here in dark box
% \end{qq}
 AFTER preamble.tex
%
% This provides:
% - Dark background with light text
% - Material Design color palette
% - Ocean blue gradient colors (p1-p9)
% - Enhanced TikZ/PGFPlots for dark backgrounds
% - qq tcolorbox environment
% - tkz-euclide for geometric constructions

% ===========================================
% DARK MODE PACKAGE
% ===========================================
\usepackage{darkmode}
\enabledarkmode

% ===========================================
% ADDITIONAL PACKAGES
% ===========================================
\usepackage{changepage}
\usepackage{ulem}
\usepackage{colortbl}
\usepackage{tkz-euclide}
\usepackage{capt-of}

% ===========================================
% EXTENDED TIKZ LIBRARIES
% ===========================================
\usetikzlibrary{patterns,3d,backgrounds}
\usepgfplotslibrary{groupplots}

% Upgrade pgfplots compatibility
\pgfplotsset{compat=1.18}

% ===========================================
% OCEAN BLUE GRADIENT (p1-p9)
% ===========================================
\definecolor{p1}{HTML}{caf0f8}
\definecolor{p2}{HTML}{ade8f4}
\definecolor{p3}{HTML}{90e0ef}
\definecolor{p4}{HTML}{48cae4}
\definecolor{p5}{HTML}{00b4d8}
\definecolor{p6}{HTML}{0096c7}
\definecolor{p7}{HTML}{0077b6}
\definecolor{p8}{HTML}{023e8a}
\definecolor{p9}{HTML}{03045e}

% Dark background color
\definecolor{pag}{HTML}{293133}

% ===========================================
% MATERIAL DESIGN COLORS
% ===========================================
% Note: myred, myyellow already defined in main preamble
% These are additional/alternate versions
\definecolor{matred}{HTML}{f44336}
\definecolor{matpink}{HTML}{e81e63}
\definecolor{matpurple}{HTML}{9c27b0}
\definecolor{matdeeppurple}{HTML}{673ab7}
\definecolor{matindigo}{HTML}{3f51b5}
\definecolor{matblue}{HTML}{2196f3}
\definecolor{matlightblue}{HTML}{03a9f4}
\definecolor{matcyan}{HTML}{00bcd4}
\definecolor{matteal}{HTML}{009688}
\definecolor{matgreen}{HTML}{4caf50}
\definecolor{matlightgreen}{HTML}{8bc34a}
\definecolor{matlime}{HTML}{cddc39}
\definecolor{matyellow}{HTML}{ffeb3b}
\definecolor{matamber}{HTML}{ffc107}
\definecolor{matorange}{HTML}{ff9800}
\definecolor{matdeeporange}{HTML}{ff5722}
\definecolor{matbrown}{HTML}{795548}
\definecolor{matgray}{HTML}{9e9e9e}
\definecolor{matbluegray}{HTML}{607d8b}

% ===========================================
% COLORLET FOR TIKZ
% ===========================================
\colorlet{xcol}{blue!60!black}

% ===========================================
% DARK MODE TCOLORBOX
% ===========================================
\newtcolorbox{qq}[2][]{%
    boxrule=0.75pt,
    sharp corners,
    colframe=white,
    colback=pag,
    coltext=white,
    #1,
}

% Dark definition box
\newtcolorbox{darkdfn}[2][]{%
    enhanced,
    boxrule=1pt,
    sharp corners,
    colframe=p5,
    colback=pag,
    coltext=white,
    coltitle=white,
    fonttitle=\bfseries,
    title=#2,
    #1,
}

% Dark theorem box
\newtcolorbox{darkthm}[2][]{%
    enhanced,
    boxrule=1pt,
    sharp corners,
    colframe=matpurple,
    colback=pag,
    coltext=white,
    coltitle=white,
    fonttitle=\bfseries,
    title=#2,
    #1,
}

% ===========================================
% TIKZ 3D FIX
% ===========================================
% Small fix for canvas is xy plane at z
% https://tex.stackexchange.com/a/48776/121799
\makeatletter
\tikzoption{canvas is xy plane at z}[]{%
    \def\tikz@plane@origin{\pgfpointxyz{0}{0}{#1}}%
    \def\tikz@plane@x{\pgfpointxyz{1}{0}{#1}}%
    \def\tikz@plane@y{\pgfpointxyz{0}{1}{#1}}%
    \tikz@canvas@is@plane}
\makeatother

% ===========================================
% CONDITIONAL LINE TIKZSET
% ===========================================
\tikzset{
    conditional line/.style 2 args={
        /utils/exec={\pgfmathsetmacro{\xcoordA}{#1}},
        /utils/exec={\pgfmathsetmacro{\xcoordB}{#2}},
        insert path={
            \ifdim\xcoordA pt>\xcoordB pt
            (\xcoordA,0) -- (\xcoordB,0)
            \fi
        },
    },
}

% ===========================================
% PGFPLOTS DARK MODE DEFAULTS
% ===========================================
\pgfplotsset{
    dark axis/.style={
        axis line style={white},
        tick style={white},
        label style={white},
        grid style={pag!70},
        every axis label/.append style={white},
        every tick label/.append style={white},
    },
}

% ===========================================
% TIKZ EXTERNALIZATION (for fast compilation)
% ===========================================
% This creates .md5 cache files for figures
% Requires: pdflatex -shell-escape
%
% To enable, add AFTER loading this file:
%   \enabletikzexternalize
%
\newcommand{\enabletikzexternalize}{%
  \usetikzlibrary{external}%
  \tikzexternalize[prefix=figures/]%
}

% ===========================================
% CONVENIENT SHORTCUTS FOR DARK PLOTS
% ===========================================
\newcommand{\darkplot}[1]{%
    \begin{tikzpicture}
    \begin{axis}[
        dark axis,
        grid=both,
        major grid style={line width=.2pt,draw=pag!70},
        #1
    ]
}


% ===========================================
% QUICK GEOMETRY MACROS (tkz-euclide)
% ===========================================
% Draw a point: \point{name}{x}{y}
\newcommand{\gpoint}[3]{\tkzDefPoint(#2,#3){#1}}

% Draw line through points: \gline{A}{B}
\newcommand{\gline}[2]{\tkzDrawLine[color=p4](#1,#2)}

% Draw circle: \gcircle{center}{radius}
\newcommand{\gcircle}[2]{\tkzDrawCircle[color=p5](#1,#2)}

% Mark angle: \gangle{A}{B}{C}
\newcommand{\gangle}[3]{\tkzMarkAngle[color=p3,size=0.5](#1,#2,#3)}

% ===========================================
% USAGE EXAMPLE
% ===========================================
% In your document:
%
% \begin{tikzpicture}
%   \begin{axis}[dark axis, ...]
%     \addplot[p4, thick] {sin(deg(x))};
%   \end{axis}
% \end{tikzpicture}
%
% Or with qq box:
% \begin{qq}{My Title}
%   Content here in dark box
% \end{qq}


\course{Studies-in-Algebra-and-Group-Theory}

\begin{document}

\lecture{6}{Kernals, images, and more}

\section{Kernel and Image}

Every linear map carries two intrinsic subspaces that reveal its fundamental structure: the \emph{kernel} (what gets annihilated) and the \emph{image} (what gets hit). These are the linear-algebraic analogues of concepts we've seen before: the kernel of a group homomorphism and the image of any function. The intermingling betwixt (what a fun word) kernel and image is governed by the \textbf{\emph{Rank-Nullity Theorem}}.

\nt{From the categorical viewpoint, kernel and image are two ways to ``measure'' a morphism. The kernel measures the \emph{failure of injectivity}---how much information the map destroys. The image measures the \emph{failure of surjectivity}---how much of the codomain the map misses. The Rank-Nullity Theorem says these two failures are complementary: the more a linear map kills, the less it can produce, and vice versa. A linear map \(\varphi : V \to W\) is completely determined by: (1) which subspace it kills (the kernel), and (2) which subspace it hits (the image). Once you know these, there's essentially only one way to fill in the rest---the map restricted to any complement of the kernel is an isomorphism onto the image. This is different from nonlinear maps, where knowing the zero set tells you almost nothing about the map's behavior elsewhere.
}

\subsection{The Kernel}

\dfn{Kernel of a linear map}{Let \(\varphi : V \to W\) be a linear map. The \textbf{\emph{kernel}} of \(\varphi\) is
\[\ker(\varphi) = \{v \in V : \varphi(v) = 0_W\}.\]
This is the set of all vectors that \(\varphi\) maps to zero.
}


\nt{The terminology ``kernel'' comes from algebra (group homomorphisms), while ``null space'' is more common in applied linear algebra and matrix theory to refer to the generalized kernal (we will get to this later). In categorical language, the kernel is the \emph{equalizer} of \(\varphi\) and the zero map.
}

\ex{Examples of kernels}{
\begin{enumerate}
  \item The zero map \(0 : V \to W\) has \(\ker(0) = V\).
  \item The identity map \(\id_V : V \to V\) has \(\ker(\id_V) = \{0\}\).
  \item Let \(\varphi : \RR^3 \to \RR\) be defined by \(\varphi(x, y, z) = x + 2y + 3z\). Then \(\ker(\varphi) = \{(x, y, z) \in \RR^3 : x + 2y + 3z = 0\}\), a plane through the origin.
  \item Let \(D : \mcP(\RR) \to \mcP(\RR)\) be differentiation, \(Dp = p'\). Then \(\ker(D)\) is the set of constant polynomials.
  \item Let \(T : \mcM_{n \times n}(k) \to k\) be the trace map, \(T(A) = Tr(A)\). Then \(\ker(T)\) is the space of traceless matrices.
\end{enumerate}
}

\mprop{The kernel is a subspace}{Let \(\varphi : V \to W\) be a linear map. Then \(\ker(\varphi)\) is a subspace of \(V\).
}

\pf{Proof}{We verify the subspace criteria:
\begin{enumerate}
  \item \textbf{Contains zero:} Since \(\varphi(0_V) = 0_W\) (every linear map preserves zero), we have \(0_V \in \ker(\varphi)\).
  \item \textbf{Closed under addition:} If \(u, v \in \ker(\varphi)\), then \(\varphi(u + v) = \varphi(u) + \varphi(v) = 0 + 0 = 0\), so \(u + v \in \ker(\varphi)\).
  \item \textbf{Closed under scalar multiplication:} If \(v \in \ker(\varphi)\) and \(\lambda \in k\), then \(\varphi(\lambda v) = \lambda \varphi(v) = \lambda \cdot 0 = 0\), so \(\lambda v \in \ker(\varphi)\).
\end{enumerate}
}


\nt{Compare this with group theory: the kernel of a group homomorphism \(f : G \to H\) is a \emph{normal} subgroup of \(G\). For linear maps, every subspace is ``normal'' in the appropriate sense---vector spaces are abelian, so there's no distinction between left and right cosets.
}

The kernel provides a clean characterization of injectivity:

\thm{Injectivity via the kernel}{Let \(\varphi : V \to W\) be a linear map. Then
\[\varphi \text{ is injective} \iff \ker(\varphi) = \{0\}.\]
}

\pf{Proof}{\((\Rightarrow)\) Suppose \(\varphi\) is injective. Let \(v \in \ker(\varphi)\). Then \(\varphi(v) = 0 = \varphi(0)\). Since \(\varphi\) is injective, \(v = 0\). Thus \(\ker(\varphi) \subseteq \{0\}\), and since \(0 \in \ker(\varphi)\) always, we have \(\ker(\varphi) = \{0\}\).

\((\Leftarrow)\) Suppose \(\ker(\varphi) = \{0\}\). Let \(u, v \in V\) with \(\varphi(u) = \varphi(v)\). Then
\[\varphi(u - v) = \varphi(u) - \varphi(v) = 0,\]
so \(u - v \in \ker(\varphi) = \{0\}\). Thus \(u - v = 0\), i.e., \(u = v\). Hence \(\varphi\) is injective.
}

\cor{}{An injective linear map from \(V\) to \(W\) embeds \(V\) as a subspace of \(W\). More precisely, if \(\varphi : V \to W\) is injective, then \(V \cong \Img(\varphi) \subseteq W\).
}

\nt{Here the kernal has a role more that merely verifying injectivity---it \emph{measures the failure of injectivity} in a precise sense. Two vectors \(u, v \in V\) have the same image under \(\varphi\) if and only if \(u - v \in \ker(\varphi)\). Thus the fibers of \(\varphi\) (preimages of points) are exactly the cosets \(v + \ker(\varphi)\). The map \(\varphi\) is injective precisely when these cosets are singletons, i.e., when \(\ker(\varphi) = \{0\}\). This is why we can ``mod out by the kernel'' to obtain an injective map---the quotient \(V/\ker(\varphi)\) collapses each fiber to a point.
}

\subsection{The Image (Range)}

\dfn{Image of a linear map}{Let \(\varphi : V \to W\) be a linear map. The \textbf{\emph{image}} (or \textbf{\emph{range}}) of \(\varphi\) is
\[\Img(\varphi) = \{\varphi(v) : v \in V\} = \{w \in W : w = \varphi(v) \text{ for some } v \in V\}.\]
This is the set of all vectors that \(\varphi\) ``hits.''
}


\ex{Examples of images}{
\begin{enumerate}
  \item The zero map \(0 : V \to W\) has \(\Img(0) = \{0_W\}\).
  \item The identity map \(\id_V : V \to V\) has \(\Img(\id_V) = V\).
  \item Let \(\pi : \RR^3 \to \RR^2\) be projection \(\pi(x, y, z) = (x, y)\). Then \(\Img(\pi) = \RR^2\).
  \item Let \(D : \mcP_n(\RR) \to \mcP_{n-1}(\RR)\) be differentiation. Then \(\Img(D) = \mcP_{n-1}(\RR)\) (every polynomial of degree \(\leq n-1\) is the derivative of some polynomial of degree \(\leq n\)).
  \item Let \(A \in \mcM_{m \times n}(k)\) define \(L_A : k^n \to k^m\). Then \(\Img(L_A) = \operatorname{ColSpace}(A)\), the column space of \(A\).
\end{enumerate}
}

\mprop{The image is a subspace}{Let \(\varphi : V \to W\) be a linear map. Then \(\Img(\varphi)\) is a subspace of \(W\).
}

\pf{Proof}{We verify the subspace criteria:
\begin{enumerate}
  \item \textbf{Contains zero:} \(\varphi(0_V) = 0_W\), so \(0_W \in \Img(\varphi)\).
  \item \textbf{Closed under addition:} If \(w_1, w_2 \in \Img(\varphi)\), then \(w_1 = \varphi(v_1)\) and \(w_2 = \varphi(v_2)\) for some \(v_1, v_2 \in V\). Thus \(w_1 + w_2 = \varphi(v_1) + \varphi(v_2) = \varphi(v_1 + v_2) \in \Img(\varphi)\).
  \item \textbf{Closed under scalar multiplication:} If \(w \in \Img(\varphi)\) and \(\lambda \in k\), then \(w = \varphi(v)\) for some \(v \in V\), so \(\lambda w = \lambda \varphi(v) = \varphi(\lambda v) \in \Img(\varphi)\).
\end{enumerate}
}

\dfn{Surjectivity}{A linear map \(\varphi : V \to W\) is \textbf{\emph{surjective}} (or \textbf{\emph{onto}}) if \(\Img(\varphi) = W\), i.e., every element of \(W\) is hit by some element of \(V\).
}

\mprop{Characterization of surjectivity}{Let \(\varphi : V \to W\) be a linear map. The following are equivalent:
\begin{enumerate}
  \item \(\varphi\) is surjective.
  \item \(\Img(\varphi) = W\).
  \item If \(\mathcal{B}\) spans \(V\), then \(\varphi(\mathcal{B}) = \{\varphi(v) : v \in \mathcal{B}\}\) spans \(W\).
\end{enumerate}
}


\subsection{The Rank-Nullity Theorem}

We now come to the central result connecting kernel and image. This theorem quantifies the fundamental trade-off in linear algebra: the more a map ``kills,'' the less it can ``produce.''

\dfn{Rank and nullity}{Let \(\varphi : V \to W\) be a linear map between finite-dimensional vector spaces. The \textbf{\emph{rank}} of \(\varphi\) is
\[\rank(\varphi) = \dim(\Img(\varphi)),\]
and the \textbf{\emph{nullity}} of \(\varphi\) is
\[\nullity(\varphi) = \dim(\ker(\varphi)).\]
For a matrix \(A\), we define \(\rank(A)\) and \(\nullity(A)\) to be the rank and nullity of the corresponding linear map.
}

\thm{The Rank-Nullity Theorem (Fundamental Theorem of Linear Maps)}{Let \(\varphi : V \to W\) be a linear map with \(V\) finite-dimensional. Then \(\Img(\varphi)\) is finite-dimensional and
\[\dim(V) = \dim(\ker(\varphi)) + \dim(\Img(\varphi)).\]
Equivalently: \(\dim(V) = \nullity(\varphi) + \rank(\varphi)\).
}

\pf{Proof}{Let \(\{u_1, \ldots, u_m\}\) be a basis for \(\ker(\varphi)\). By the basis extension theorem, we can extend this to a basis
\[\{u_1, \ldots, u_m, v_1, \ldots, v_r\}\]
of \(V\), where \(\dim(V) = m + r\).

\textbf{Claim:} \(\{\varphi(v_1), \ldots, \varphi(v_r)\}\) is a basis for \(\Img(\varphi)\).

\emph{Spanning:} Let \(w \in \Img(\varphi)\), say \(w = \varphi(x)\) for some \(x \in V\). Write
\[x = a_1 u_1 + \cdots + a_m u_m + b_1 v_1 + \cdots + b_r v_r.\]
Applying \(\varphi\):
\[w = \varphi(x) = a_1 \varphi(u_1) + \cdots + a_m \varphi(u_m) + b_1 \varphi(v_1) + \cdots + b_r \varphi(v_r).\]
Since each \(u_i \in \ker(\varphi)\), we have \(\varphi(u_i) = 0\), so
\[w = b_1 \varphi(v_1) + \cdots + b_r \varphi(v_r).\]
Thus \(\{\varphi(v_1), \ldots, \varphi(v_r)\}\) spans \(\Img(\varphi)\).

\emph{Linear independence:} Suppose
\[c_1 \varphi(v_1) + \cdots + c_r \varphi(v_r) = 0.\]
Then \(\varphi(c_1 v_1 + \cdots + c_r v_r) = 0\), so \(c_1 v_1 + \cdots + c_r v_r \in \ker(\varphi)\). Since \(\{u_1, \ldots, u_m\}\) is a basis for \(\ker(\varphi)\), we can write
\[c_1 v_1 + \cdots + c_r v_r = d_1 u_1 + \cdots + d_m u_m\]
for some scalars \(d_j\). Rearranging:
\[d_1 u_1 + \cdots + d_m u_m - c_1 v_1 - \cdots - c_r v_r = 0.\]
Since \(\{u_1, \ldots, u_m, v_1, \ldots, v_r\}\) is linearly independent (it's a basis of \(V\)), all coefficients are zero. In particular, \(c_1 = \cdots = c_r = 0\).

Thus \(\{\varphi(v_1), \ldots, \varphi(v_r)\}\) is a basis for \(\Img(\varphi)\), so \(\dim(\Img(\varphi)) = r\). We conclude:
\[\dim(V) = m + r = \dim(\ker(\varphi)) + \dim(\Img(\varphi)).\]
}


\nt{The Rank-Nullity Theorem is the linear algebra version of the \textbf{\emph{First Isomorphism Theorem}} from group theory. Recall that for a group homomorphism \(f : G \to H\), we have \(G/\ker(f) \cong \Img(f)\). The linear version states \(V/\ker(\varphi) \cong \Img(\varphi)\), and ``counting dimensions'' gives \(\dim(V) - \dim(\ker(\varphi)) = \dim(\Img(\varphi))\).
}

\subsection{Consequences of the Rank-Nullity Theorem}

The Rank-Nullity Theorem has powerful implications for understanding linear maps.

\cor{Dimension constraints on injectivity}{Let \(\varphi : V \to W\) be a linear map with \(\dim(V) = n\) and \(\dim(W) = m\).
\begin{enumerate}
  \item If \(n > m\), then \(\varphi\) cannot be injective.
  \item If \(n < m\), then \(\varphi\) cannot be surjective.
\end{enumerate}
}

\pf{Proof}{(1) If \(\varphi\) were injective, then \(\dim(\ker(\varphi)) = 0\), so by Rank-Nullity:
\[\dim(\Img(\varphi)) = \dim(V) = n > m = \dim(W).\]
But \(\Img(\varphi) \subseteq W\), so \(\dim(\Img(\varphi)) \leq \dim(W) = m\). Contradiction.

(2) If \(\varphi\) were surjective, then \(\dim(\Img(\varphi)) = \dim(W) = m\), so by Rank-Nullity:
\[\dim(\ker(\varphi)) = \dim(V) - \dim(\Img(\varphi)) = n - m < 0.\]
But dimensions are non-negative. Contradiction.
}

\cor{Equal dimensions: bijectivity criteria}{Let \(\varphi : V \to W\) be a linear map with \(\dim(V) = \dim(W) = n\). The following are equivalent:
\begin{enumerate}
  \item \(\varphi\) is injective.
  \item \(\varphi\) is surjective.
  \item \(\varphi\) is bijective (an isomorphism).
  \item \(\ker(\varphi) = \{0\}\).
  \item \(\Img(\varphi) = W\).
  \item \(\rank(\varphi) = n\).
\end{enumerate}
}

\pf{Proof}{By Rank-Nullity, \(\nullity(\varphi) + \rank(\varphi) = n\). Since \(\rank(\varphi) \leq \dim(W) = n\):
\begin{itemize}
  \item \(\varphi\) injective \(\iff \nullity(\varphi) = 0 \iff \rank(\varphi) = n\).
  \item \(\varphi\) surjective \(\iff \rank(\varphi) = n \iff \nullity(\varphi) = 0\).
\end{itemize}
Thus injectivity, surjectivity, and bijectivity are equivalent when dimensions match.
}

\nt{For linear maps between equal-dimensional spaces, ``one-to-one'' and ``onto'' are the same condition. This fails for general functions and even for linear maps between infinite-dimensional spaces. Consider the \textbf{\emph{right shift}} \(R : k^\infty \to k^\infty\) defined by \(R(x_1, x_2, x_3, \ldots) = (0, x_1, x_2, \ldots)\). This is injective (kernel is trivial) but not surjective (the sequence \((1, 0, 0, \ldots)\) is not in the image). Conversely, the \textbf{\emph{left shift}} \(L(x_1, x_2, x_3, \ldots) = (x_2, x_3, \ldots)\) is surjective but not injective. In infinite dimensions, there's ``room'' for one without the other. The finite-dimensional equivalence is a manifestation of the pigeonhole principle applied to basis vectors.
}


\cor{Existence of right/left inverses}{Let \(\varphi : V \to W\) be a linear map between finite-dimensional spaces.
\begin{enumerate}
  \item \(\varphi\) has a left inverse if and only if \(\varphi\) is injective.
  \item \(\varphi\) has a right inverse if and only if \(\varphi\) is surjective.
\end{enumerate}
}

\subsection{Matrix Interpretation}

For a matrix \(A \in \mcM_{m \times n}(k)\), the associated linear map is \(L_A : k^n \to k^m\) given by \(L_A(x) = Ax\). The Rank-Nullity Theorem becomes:

\mprop{Rank-Nullity for matrices}{For \(A \in \mcM_{m \times n}(k)\):
\begin{enumerate}
  \item \(\ker(L_A) = \{x \in k^n : Ax = 0\}\) is the \emph{null space} of \(A\).
  \item \(\Img(L_A) = \{Ax : x \in k^n\}\) is the \emph{column space} of \(A\) (the span of the columns).
  \item \(\nullity(A) + \rank(A) = n\) (the number of columns).
\end{enumerate}
}

\nt{The rank of a matrix equals the number of linearly independent columns, which equals the number of linearly independent rows. The column space has dimension equal to the rank, and the null space has dimension equal to \(n - \rank(A)\).

Looking ahead (citing concepts we'll introduce in later lectures): when we study eigenvalues, the kernel will reappear in a crucial role. For a linear operator \(T : V \to V\) and scalar \(\lambda\), the \textbf{\emph{eigenspace}} \(E_\lambda = \ker(T - \lambda I)\) is precisely the kernel of the shifted operator. The dimension of this eigenspace (the \emph{geometric multiplicity} of \(\lambda\)) measures ``how degenerate'' the eigenvalue is. The interplay between algebraic and geometric multiplicities---and the failure of these to match---is what makes Jordan normal form necessary.
}

\ex{Computing kernel and image}{Let \(A = \begin{pmatrix} 1 & 2 & 3 \\ 0 & 1 & 2 \\ 1 & 3 & 5 \end{pmatrix}\) over \(\RR\).

\textbf{Finding \(\ker(A)\):} Solve \(Ax = 0\). Row reduce:
\[\begin{pmatrix} 1 & 2 & 3 \\ 0 & 1 & 2 \\ 1 & 3 & 5 \end{pmatrix} \xrightarrow{R_3 - R_1} \begin{pmatrix} 1 & 2 & 3 \\ 0 & 1 & 2 \\ 0 & 1 & 2 \end{pmatrix} \xrightarrow{R_3 - R_2} \begin{pmatrix} 1 & 2 & 3 \\ 0 & 1 & 2 \\ 0 & 0 & 0 \end{pmatrix}.\]
Back-substituting: \(x_2 = -2x_3\), \(x_1 = -2x_2 - 3x_3 = 4x_3 - 3x_3 = x_3\). So
\[\ker(A) = \operatorname{Span}\{(1, -2, 1)^T\}, \quad \nullity(A) = 1.\]

\textbf{Finding \(\Img(A)\):} The non-zero rows in RREF correspond to pivot columns. Columns 1 and 2 are pivot columns, so
\[\Img(A) = \operatorname{Span}\left\{\begin{pmatrix} 1 \\ 0 \\ 1 \end{pmatrix}, \begin{pmatrix} 2 \\ 1 \\ 3 \end{pmatrix}\right\}, \quad \rank(A) = 2.\]

\textbf{Verification:} \(\nullity(A) + \rank(A) = 1 + 2 = 3 = n\). \checkmark
}


\qs{}{
\begin{enumerate}
  \item Let \(D : \mcP_3(\RR) \to \mcP_3(\RR)\) be the differentiation operator. Find \(\ker(D)\), \(\Img(D)\), \(\nullity(D)\), and \(\rank(D)\). Verify the Rank-Nullity Theorem.
  \item Let \(T : \RR^4 \to \RR^3\) be defined by \(T(x_1, x_2, x_3, x_4) = (x_1 - x_2, x_2 - x_3, x_3 - x_4)\). Determine whether \(T\) is injective, surjective, or neither.
  \item Prove that if \(\varphi : V \to V\) is a linear operator on a finite-dimensional space with \(\varphi^2 = \varphi\), then \(V = \ker(\varphi) \oplus \Img(\varphi)\).
  \item Let \(V\) and \(W\) be finite-dimensional vector spaces with \(\dim(V) = \dim(W)\). Prove that a linear map \(\varphi : V \to W\) is injective if and only if for every linearly independent list in \(V\), its image under \(\varphi\) is linearly independent in \(W\).
\end{enumerate}
}

\begin{solution}
\textbf{(1)} The space \(\mcP_3(\RR)\) has basis \(\{1, x, x^2, x^3\}\), so \(\dim(\mcP_3(\RR)) = 4\).

\textbf{Kernel:} \(Dp = 0\) if and only if \(p\) is a constant polynomial. So \(\ker(D) = \operatorname{Span}\{1\}\), which are the constant functions. Thus \(\nullity(D) = 1\).

\textbf{Image:} Differentiating: \(D(1) = 0\), \(D(x) = 1\), \(D(x^2) = 2x\), \(D(x^3) = 3x^2\). The non-zero outputs span \(\operatorname{Span}\{1, x, x^2\} = \mcP_2(\RR)\). Thus \(\Img(D) = \mcP_2(\RR)\) and \(\rank(D) = 3\).

\textbf{Verification:} \(\nullity(D) + \rank(D) = 1 + 3 = 4 = \dim(\mcP_3(\RR))\). \checkmark

\textbf{(2)} \textbf{Injectivity:} Find \(\ker(T)\). Solve \(T(x_1, x_2, x_3, x_4) = (0, 0, 0)\):
\[x_1 - x_2 = 0, \quad x_2 - x_3 = 0, \quad x_3 - x_4 = 0.\]
This gives \(x_1 = x_2 = x_3 = x_4\). So \(\ker(T) = \{(t, t, t, t) : t \in \RR\} = \operatorname{Span}\{(1, 1, 1, 1)\}\).

Since \(\ker(T) \neq \{0\}\), \(T\) is \textbf{not injective}.

\textbf{Surjectivity:} We have \(\nullity(T) = 1\), so \(\rank(T) = 4 - 1 = 3 = \dim(\RR^3)\). Thus \(\Img(T) = \RR^3\), and \(T\) is \textbf{surjective}.

\textbf{(3)} We show \(V = \ker(\varphi) + \Img(\varphi)\) and \(\ker(\varphi) \cap \Img(\varphi) = \{0\}\).

\emph{Sum equals \(V\):} Let \(v \in V\). Write \(v = (v - \varphi(v)) + \varphi(v)\). Note that \(\varphi(v) \in \Img(\varphi)\). Also:
\[\varphi(v - \varphi(v)) = \varphi(v) - \varphi^2(v) = \varphi(v) - \varphi(v) = 0,\]
so \(v - \varphi(v) \in \ker(\varphi)\). Thus \(v \in \ker(\varphi) + \Img(\varphi)\).

\emph{Intersection is trivial:} Let \(w \in \ker(\varphi) \cap \Img(\varphi)\). Then \(\varphi(w) = 0\) and \(w = \varphi(u)\) for some \(u \in V\). So:
\[w = \varphi(u) = \varphi^2(u) = \varphi(\varphi(u)) = \varphi(w) = 0.\]

Thus \(V = \ker(\varphi) \oplus \Img(\varphi)\).

\textbf{(4)} (\(\Rightarrow\)) Suppose \(\varphi\) is injective and \(\{v_1, \ldots, v_m\}\) is linearly independent. Suppose \(\sum c_i \varphi(v_i) = 0\). Then \(\varphi(\sum c_i v_i) = 0\), so \(\sum c_i v_i \in \ker(\varphi) = \{0\}\). Thus \(\sum c_i v_i = 0\), and by linear independence of the \(v_i\), all \(c_i = 0\). Hence \(\{\varphi(v_1), \ldots, \varphi(v_m)\}\) is linearly independent.

(\(\Leftarrow\)) Suppose \(\varphi\) maps every linearly independent set to a linearly independent set. If \(v \neq 0\), then \(\{v\}\) is linearly independent, so \(\{\varphi(v)\}\) is linearly independent, meaning \(\varphi(v) \neq 0\). Thus \(\ker(\varphi) = \{0\}\), so \(\varphi\) is injective.
\end{solution}


\section{Duality}

The \emph{dual space} of a vector space \(V\) consists of all linear maps from \(V\) to the base field \(k\)---these are called \emph{linear functionals}. While this might seem like a minor variation on Hom-spaces, duality reveals useful structure: every basis has a ``dual basis,'' every linear map has a ``dual map,'' and the matrix of the dual map is the transpose of the original. In more advanced contexts (differential geometry, functional analysis, representation theory), duality becomes indispensable.

\nt{The dual space construction is the first example of a \textbf{\emph{contravariant functor}} you'll encounter. When we pass from \(V\) to \(V'\), linear maps \(T : V \to W\) don't induce maps \(V' \to W'\)---instead, they induce maps \(W' \to V'\), going the ``wrong way.'' This reversal of arrows is the hallmark of contravariance. In category theory, the dual space functor \((-)' : \mathbf{Vect}_k^{\mathrm{op}} \to \mathbf{Vect}_k\) is representable: \(V' = \Hom(V, k)\) is ``the functor represented by \(k\).''
}

\subsection{Linear Functionals and the Dual Space}

\dfn{Linear functional}{A \textbf{\emph{linear functional}} on a vector space \(V\) over a field \(k\) is a linear map \(\ell : V \to k\). In other words, a linear functional is an element of \(\Hom(V, k)\).
}

\ex{Examples of linear functionals}{
\begin{enumerate}
  \item \textbf{Coordinate projections:} On \(k^n\), define \(\pi_j(x_1, \ldots, x_n) = x_j\). Each \(\pi_j\) is a linear functional.
  \item \textbf{Weighted sums:} Fix \((c_1, \ldots, c_n) \in k^n\). Define \(\ell(x_1, \ldots, x_n) = c_1 x_1 + \cdots + c_n x_n\). This is a linear functional on \(k^n\), and \emph{every} linear functional on \(k^n\) has this form.
  \item \textbf{Evaluation:} On \(\mcP(\RR)\), define \(\ev_a(p) = p(a)\) for fixed \(a \in \RR\). This is a linear functional.
  \item \textbf{Integration:} On \(\mcP(\RR)\), define \(\ell(p) = \int_0^1 p(x) \, dx\). Linearity of the integral makes this a linear functional.
  \item \textbf{Trace:} On \(\mcM_{n \times n}(k)\), the trace \(Tr(A) = \sum_{i=1}^n a_{ii}\) is a linear functional.
\end{enumerate}
}

\dfn{Dual space}{The \textbf{\emph{dual space}} of \(V\), denoted \(V^*\) or \(V'\), is the vector space of all linear functionals on \(V\):
\[V^* = \Hom(V, k).\]
This is a vector space under pointwise addition and scalar multiplication: \((\ell_1 + \ell_2)(v) = \ell_1(v) + \ell_2(v)\) and \((\lambda \ell)(v) = \lambda \cdot \ell(v)\).
}

\nt{We use \(V^*\) notation (common in differential geometry and physics) interchangeably with \(V'\) (common in algebra). Both mean the same thing. Some texts reserve \(V^*\) for the continuous dual in functional analysis, but for finite-dimensional spaces this distinction is irrelevant.
}

\mprop{Dimension of the dual space}{If \(V\) is finite-dimensional with \(\dim(V) = n\), then \(\dim(V^*) = n\). Thus \(V^* \cong V\) as vector spaces (non-canonically).
}

\pf{Proof}{By our earlier result on Hom-spaces:
\[\dim(V^*) = \dim(\Hom(V, k)) = \dim(V) \cdot \dim(k) = n \cdot 1 = n.\]
}

\nt{The isomorphism \(V \cong V^*\) is \emph{non-canonical}---it depends on a choice of basis. There's no ``natural'' way to identify vectors with linear functionals without extra structure. However, the \emph{double dual} \(V^{**} = (V^*)^*\) \emph{does} have a canonical isomorphism with \(V\), given by evaluation: \(v \mapsto (\ell \mapsto \ell(v))\).
}

\subsection{The Dual Basis}

Given a basis of \(V\), there's a natural associated basis of \(V^*\).

\dfn{Dual basis}{Let \(\mathcal{B} = \{v_1, \ldots, v_n\}\) be a basis of \(V\). The \textbf{\emph{dual basis}} is the list \(\mathcal{B}^* = \{\varphi_1, \ldots, \varphi_n\}\) of linear functionals in \(V^*\) defined by
\[\varphi_j(v_i) = \delta_{ij} = \begin{cases} 1 & \text{if } i = j, \\ 0 & \text{if } i \neq j. \end{cases}\]
In other words, \(\varphi_j\) ``picks out'' the \(j\)-th coordinate: if \(v = c_1 v_1 + \cdots + c_n v_n\), then \(\varphi_j(v) = c_j\).
}

\mprop{Dual basis is a basis}{The dual basis \(\{\varphi_1, \ldots, \varphi_n\}\) is indeed a basis of \(V^*\).
}

\pf{Proof}{\textbf{Linear independence:} Suppose \(a_1 \varphi_1 + \cdots + a_n \varphi_n = 0\) (the zero functional). Evaluating at \(v_k\):
\[0 = (a_1 \varphi_1 + \cdots + a_n \varphi_n)(v_k) = a_1 \varphi_1(v_k) + \cdots + a_n \varphi_n(v_k) = a_k.\]
So all \(a_k = 0\).

\textbf{Spanning:} Since we have \(n\) linearly independent vectors in \(V^*\) and \(\dim(V^*) = n\), they form a basis.
}

\mprop{Dual basis extracts coordinates}{Let \(\{v_1, \ldots, v_n\}\) be a basis of \(V\) with dual basis \(\{\varphi_1, \ldots, \varphi_n\}\). For any \(v \in V\):
\[v = \varphi_1(v) v_1 + \varphi_2(v) v_2 + \cdots + \varphi_n(v) v_n.\]
That is, \(\varphi_j(v)\) is the \(j\)-th coordinate of \(v\) with respect to the basis \(\{v_i\}\).
}

\pf{Proof}{Write \(v = c_1 v_1 + \cdots + c_n v_n\). Applying \(\varphi_j\):
\[\varphi_j(v) = c_1 \varphi_j(v_1) + \cdots + c_n \varphi_j(v_n) = c_j.\]
Substituting back: \(v = \varphi_1(v) v_1 + \cdots + \varphi_n(v) v_n\).
}

\ex{Dual basis of the standard basis}{Let \(\{e_1, \ldots, e_n\}\) be the standard basis of \(k^n\). The dual basis consists of the coordinate projections:
\[\varphi_j(x_1, \ldots, x_n) = x_j.\]
These are the ``row vectors'' \(\varphi_j = (0, \ldots, 0, 1, 0, \ldots, 0)\) (with 1 in position \(j\)).
}


\subsection{The Dual Map}

Just as a linear map \(T : V \to W\) induces a map on Hom-spaces, it induces a map on dual spaces---but in the \emph{opposite} direction.

\dfn{Dual map (transpose)}{Let \(T : V \to W\) be a linear map. The \textbf{\emph{dual map}} (or \textbf{\emph{transpose}}) is the linear map \(T^* : W^* \to V^*\) defined by
\[T^*(\psi) = \psi \circ T\]
for each \(\psi \in W^*\). In other words, \(T^*(\psi)\) is the linear functional on \(V\) given by \(v \mapsto \psi(T(v))\).
}

\nt{The dual map goes ``backwards'': \(T : V \to W\) yields \(T^* : W^* \to V^*\). This is because we're precomposing with \(T\). If you think of linear functionals as ``measurements,'' then \(T^*(\psi)\) is the measurement you get by first applying \(T\), then measuring with \(\psi\). The reversal of direction is the essence of contravariance.
}

\mprop{The dual map is linear}{For \(T : V \to W\), the dual map \(T^* : W^* \to V^*\) is indeed a linear map.
}

\pf{Proof}{For \(\psi_1, \psi_2 \in W^*\) and \(\lambda \in k\):
\begin{align*}
T^*(\psi_1 + \psi_2) &= (\psi_1 + \psi_2) \circ T = \psi_1 \circ T + \psi_2 \circ T = T^*(\psi_1) + T^*(\psi_2), \\
T^*(\lambda \psi) &= (\lambda \psi) \circ T = \lambda(\psi \circ T) = \lambda T^*(\psi).
\end{align*}
}

\mprop{Algebraic properties of the dual map}{Let \(V, W, U\) be vector spaces over \(k\).
\begin{enumerate}
  \item For \(S, T : V \to W\), \((S + T)^* = S^* + T^*\).
  \item For \(T : V \to W\) and \(\lambda \in k\), \((\lambda T)^* = \lambda T^*\).
  \item For \(S : V \to W\) and \(T : W \to U\), \((T \circ S)^* = S^* \circ T^*\) (order reverses!).
  \item \((\id_V)^* = \id_{V^*}\).
\end{enumerate}
}

\pf{Proof}{(3) is the key property. For \(\psi \in U^*\):
\[(T \circ S)^*(\psi) = \psi \circ (T \circ S) = (\psi \circ T) \circ S = S^*(\psi \circ T) = S^*(T^*(\psi)) = (S^* \circ T^*)(\psi).\]
The other properties are similar.
}

\nt{Property (3) says the dual is a \emph{contravariant functor}: it reverses the order of composition. Compare with the transpose of matrices: \((AB)^T = B^T A^T\). This is not a coincidence---the matrix of the dual map \emph{is} the transpose!
}

\subsection{Matrix of the Dual Map}

The connection between dual maps and matrix transposes clarifies why the transpose operation is natural.

\thm{Matrix of the dual map is the transpose}{Let \(T : V \to W\) be a linear map between finite-dimensional spaces. Fix bases \(\mathcal{B} = \{v_1, \ldots, v_n\}\) of \(V\) and \(\mathcal{C} = \{w_1, \ldots, w_m\}\) of \(W\), with dual bases \(\mathcal{B}^* = \{\varphi_1, \ldots, \varphi_n\}\) and \(\mathcal{C}^* = \{\psi_1, \ldots, \psi_m\}\). If 
\[A = \mcM(T; \mathcal{B}, \mathcal{C}) \in \mcM_{m \times n}(k)\]
is the matrix of \(T\), then the matrix of \(T^*\) with respect to the dual bases is the transpose:
\[\mcM(T^*; \mathcal{C}^*, \mathcal{B}^*) = A^T \in \mcM_{n \times m}(k).\]
}

\pf{Proof}{We need to show that \(T^*(\psi_j) = \sum_{i=1}^n A_{ji} \varphi_i\) for each \(j\).

By definition of the matrix \(A\), we have \(T(v_k) = \sum_{i=1}^m A_{ik} w_i\). Now compute:
\[T^*(\psi_j)(v_k) = \psi_j(T(v_k)) = \psi_j\left(\sum_{i=1}^m A_{ik} w_i\right) = \sum_{i=1}^m A_{ik} \psi_j(w_i) = A_{jk}.\]

On the other hand, if \(T^*(\psi_j) = \sum_{i=1}^n c_i \varphi_i\), then
\[T^*(\psi_j)(v_k) = \sum_{i=1}^n c_i \varphi_i(v_k) = c_k.\]

Comparing: \(c_k = A_{jk}\). Thus \(T^*(\psi_j) = \sum_{k=1}^n A_{jk} \varphi_k\), which means the \(j\)-th column of \(\mcM(T^*)\) is the \(j\)-th row of \(A\). Hence \(\mcM(T^*) = A^T\).
}


\nt{This theorem explains why the transpose is natural: it's not just a matrix operation, it's the coordinate representation of the dual map. The reversal of dimensions (\(m \times n\) becomes \(n \times m\)) reflects the contravariance of the dual construction.
}

\cor{Rank of transpose equals rank}{For any matrix \(A\), \(\rank(A^T) = \rank(A)\).
}

\pf{Proof}{The rank of \(A\) is \(\dim(\Img(T))\) where \(T\) is the associated linear map. By a result we'll see shortly, \(\dim(\Img(T^*)) = \dim(\Img(T))\). Since the matrix of \(T^*\) is \(A^T\), we get \(\rank(A^T) = \rank(A)\).
}

\subsection{Kernel and Image of the Dual Map}

The dual map interchanges the roles of kernel and image in a precise way.

\mprop{Relationships between \(T\) and \(T^*\)}{Let \(T : V \to W\) be a linear map between finite-dimensional spaces.
\begin{enumerate}
  \item \(\ker(T^*) = (\Img(T))^\circ\), the \emph{annihilator} of the image.
  \item \(\Img(T^*) = (\ker(T))^\circ\), the annihilator of the kernel.
  \item \(\dim(\ker(T^*)) = \dim(W) - \rank(T)\).
  \item \(\dim(\Img(T^*)) = \rank(T)\).
\end{enumerate}
Here, for a subspace \(U \subseteq V\), the \textbf{\emph{annihilator}} is \(U^\circ = \{\ell \in V^* : \ell(u) = 0 \text{ for all } u \in U\}\).
}

\nt{The annihilator \(U^\circ\) consists of all linear functionals that vanish on \(U\). If \(\dim(V) = n\) and \(\dim(U) = k\), then \(\dim(U^\circ) = n - k\). The annihilator provides a duality between subspaces of \(V\) and subspaces of \(V^*\): bigger subspaces have smaller annihilators.
}

\subsection{The Double Dual and Natural Isomorphisms}

We've seen that \(V \cong V^*\) for finite-dimensional \(V\), but this isomorphism required a choice of basis. The double dual \(V^{**} = (V^*)^*\) enjoys a more canonical relationship with \(V\).

\dfn{Evaluation map}{The \textbf{evaluation map} is \(\ev : V \to V^{**}\) defined by
\[\ev(v)(\ell) = \ell(v)\]
for \(v \in V\) and \(\ell \in V^*\). That is, \(\ev(v)\) is the linear functional on \(V^*\) that ``evaluates at \(v\).''}

\mprop{Properties of the evaluation map}{The evaluation map \(\ev : V \to V^{**}\) satisfies:
\begin{enumerate}
  \item \(\ev\) is linear.
  \item \(\ev\) is injective.
  \item If \(V\) is finite-dimensional, \(\ev\) is an isomorphism.
\end{enumerate}
}

\pf{Proof}{\textbf{(1)} For \(v, w \in V\), \(\lambda \in k\), and \(\ell \in V^*\):
\[\ev(v + \lambda w)(\ell) = \ell(v + \lambda w) = \ell(v) + \lambda \ell(w) = \ev(v)(\ell) + \lambda \ev(w)(\ell).\]

\textbf{(2)} Suppose \(\ev(v) = 0\), meaning \(\ell(v) = 0\) for all \(\ell \in V^*\). If \(v \neq 0\), extend \(\{v\}\) to a basis and let \(\ell\) be the dual basis element with \(\ell(v) = 1\). Contradiction. So \(v = 0\).

\textbf{(3)} For finite-dimensional \(V\): \(\dim(V^{**}) = \dim(V^*) = \dim(V)\). Since \(\ev\) is injective, it's an isomorphism.
}

\nt{The crucial point: \(\ev\) is \emph{natural}---it doesn't depend on any choices. Given bases \(\{e_1, \ldots, e_n\}\) of \(V\), dual basis \(\{e_1^*, \ldots, e_n^*\}\) of \(V^*\), and double dual basis \(\{e_1^{**}, \ldots, e_n^{**}\}\) of \(V^{**}\), one can verify that \(\ev(e_i) = e_i^{**}\). The evaluation map sends each basis element to its ``double dual counterpart.''}

\nt{For infinite-dimensional \(V\), the situation is strikingly different. If \(V\) has a basis \(\{e_i\}_{i \in I}\) indexed by an infinite set \(I\), then every element of \(V\) is a \emph{finite} linear combination \(\sum_{i \in I} x_i e_i\) (with only finitely many \(x_i \neq 0\)). However, elements of \(V^*\) can be arbitrary: given any function \(a : I \to k\), the map \(\ell_a(\sum x_i e_i) = \sum x_i a_i\) is a well-defined linear functional. Thus \(V^* \cong k^I = \prod_{i \in I} k\), the full product, which is much larger than the direct sum \(V \cong \bigoplus_{i \in I} k\). The evaluation map \(\ev : V \to V^{**}\) is still injective, but far from surjective.}

\qs{}{
\begin{enumerate}
  \item Let \(V = \RR^3\) with standard basis \(\{e_1, e_2, e_3\}\). Write down the dual basis explicitly as functions \(\RR^3 \to \RR\).
  \item Let \(T : \RR^3 \to \RR^2\) be defined by \(T(x, y, z) = (x + y, y + z)\). Find the matrix of \(T\) with respect to the standard bases, then find the matrix of \(T^*\) with respect to the dual bases.
  \item Prove that if \(T : V \to W\) is surjective, then \(T^* : W^* \to V^*\) is injective.
  \item Prove that if \(V\) is finite-dimensional, the map \(\Phi : V \to V^{**}\) defined by \(\Phi(v)(\ell) = \ell(v)\) is an isomorphism.
\end{enumerate}
}

\begin{solution}
\textbf{(1)} The dual basis consists of the coordinate projections:
\[\varphi_1(x, y, z) = x, \quad \varphi_2(x, y, z) = y, \quad \varphi_3(x, y, z) = z.\]
These are ``row vectors'' \((1, 0, 0)\), \((0, 1, 0)\), \((0, 0, 1)\).

\textbf{(2)} The matrix of \(T\) is:
\[A = \mcM(T) = \begin{pmatrix} 1 & 1 & 0 \\ 0 & 1 & 1 \end{pmatrix}.\]
The matrix of \(T^*\) is the transpose:
\[A^T = \mcM(T^*) = \begin{pmatrix} 1 & 0 \\ 1 & 1 \\ 0 & 1 \end{pmatrix}.\]
Explicitly, if \(\psi_1, \psi_2\) are the dual basis of \(\RR^2\) and \(\varphi_1, \varphi_2, \varphi_3\) are the dual basis of \(\RR^3\):
\[T^*(\psi_1) = \varphi_1 + \varphi_2, \quad T^*(\psi_2) = \varphi_2 + \varphi_3.\]

\textbf{(3)} Suppose \(T\) is surjective and \(T^*(\psi) = 0\), i.e., \(\psi \circ T = 0\). For any \(w \in W\), since \(T\) is surjective, \(w = T(v)\) for some \(v \in V\). Then:
\[\psi(w) = \psi(T(v)) = (\psi \circ T)(v) = 0.\]
Since this holds for all \(w \in W\), we have \(\psi = 0\). Thus \(\ker(T^*) = \{0\}\), so \(T^*\) is injective.

\textbf{(4)} \textbf{Well-defined and linear:} For \(v \in V\), \(\Phi(v) : V^* \to k\) is defined by \(\Phi(v)(\ell) = \ell(v)\). This is linear in \(\ell\) (check!), so \(\Phi(v) \in V^{**}\). Moreover, \(\Phi\) is linear: \(\Phi(v_1 + v_2)(\ell) = \ell(v_1 + v_2) = \ell(v_1) + \ell(v_2) = \Phi(v_1)(\ell) + \Phi(v_2)(\ell)\).

\textbf{Injective:} Suppose \(\Phi(v) = 0\), meaning \(\ell(v) = 0\) for all \(\ell \in V^*\). If \(v \neq 0\), extend \(\{v\}\) to a basis \(\{v, v_2, \ldots, v_n\}\). The dual basis element \(\varphi_1\) satisfies \(\varphi_1(v) = 1 \neq 0\), contradiction. So \(v = 0\).

\textbf{Surjective:} Since \(\dim(V) = \dim(V^*) = \dim(V^{**})\), injectivity implies surjectivity.
\end{solution}






\end{document}
